\section{Related Work}
\label{sec:relwork}

\paragraph{AST proof rules.}
Bounded variants establish AST through a uniformly positive probability
of descent in a bounded natural-valued range~\cite{DBLP:series/mcs/McIverM05}.
McIver et al.~\cite{10.1145/3158121} give a source-level rule, including
demonic choice, in which a single real-valued function satisfies both
supermartingale and probabilistic progress obligations. Majumdar and
Sathiyanarayana~\cite{DBLP:journals/pacmpl/MajumdarS25} separate these roles
at pCFG level and prove relative completeness in arithmetic, also
establishing completeness of the earlier single-function rule.
Our contribution concerns the representation of certificates and their
transfer to whole iterations, rather than greater semantic proving power.
Separating the roles allows a resetting variant without requiring that
particular variant to be a supermartingale; it does not show that every
single-function certificate must fail.

The program-level transfer originates in Seidel's master's
thesis~\cite{seidel2026thesis}, whose Chapter~6 presents the two-certificate
rule and translations between program and pCFG certificates. The present
paper develops that foundation with explicit normalization and boundedness
arguments, sampler case studies, and a reproducible synthesis prototype.

\paragraph{Distribution transformers and program logics.}
Kozen's semantics~\cite{DBLP:conf/focs/Kozen79,DBLP:conf/stoc/Kozen83}
provides the denotational foundation for our post-distribution
transformers. Expectation-transformer reasoning~\cite{DBLP:series/mcs/McIverM05}
and distributional program logics, including den Hartog and de
Vink~\cite{DBLP:journals/ijfcs/HartogV02} and
Ellora~\cite{DBLP:conf/esop/BartheEGGHS18}, support broader probabilistic
correctness properties. Our AST rule complements such reasoning:
distributional invariants can carry output information, while termination
uses their support and two state functions. Thus, the rule does not
replace distributional output-correctness arguments.

\paragraph{Certificate and invariant synthesis.}
Constraint-based synthesis already supports martingales and ranking
supermartingales~\cite{ChakarovSankaranarayanan2013}. Batz et
al.~\cite{BatzEtAl2023} use counterexample-guided synthesis of piecewise
linear quantitative invariants at source level, including for
infinite-state programs. Their objectives include expected-outcome bounds
and finite expected runtime. We instantiate this search pattern with two AST certificates and the
sufficient coupling $U\leq V$. The random
walk illustrates qualitative termination with infinite expected runtime.
Neither CEGIS nor affine or piecewise affine templates are new here.
The contribution is their explicit instantiation for these two
iteration-level certificate roles. Invariants, transition summaries,
feature bases, and partitions remain supplied inputs; bounded search is
incomplete even though successful candidates are checked universally
over the encoded states.

Haase et al.~\cite{haase2026generatingfunctionsmeetoccupation} synthesize
occupation invariants represented by generating functions. Such
invariants describe expected visit counts and support exact inference,
with positive AST obtainable under suitable finiteness conditions.
Their unknowns and verification conditions therefore differ from our
state-function certificates for qualitative AST. Template instantiation
is common to both approaches, but their results are not interchangeable.

\paragraph{Exact samplers and their verification.}
The algorithms studied here are established constructions:
Lumbroso's FDR~\cite{DBLP:journals/corr/abs-1304-1916}, Saad et al.'s
FLDR~\cite{DBLP:conf/aistats/SaadFRM20}, and Han and Hoshi's interval
algorithm~\cite{HanHoshi1997}. Zilken et
al.~\cite{zilken2026verifyingsamplingalgorithmsdistributional} already
verify partial and total correctness of FDR and FLDR using
distributional invariants. We reuse those invariants and revisit AST
through the two-certificate rule. Han--Hoshi supplies a different
termination argument, rather than a new output-correctness proof.
Bagnall et al.'s Zar compiler~\cite{DBLP:journals/pacmpl/Bagnall0023}
verifies compilation of probabilistic programs into random-bit samplers
in Coq. Our handwritten proofs and solver-backed prototype do not
provide that level of mechanization.

\paragraph{Nondeterminism.}
Both AST rules above support demonic nondeterminism. Our current
transformer assigns one post-distribution to each input. Extending the
rule to universal scheduler quantification requires corresponding
semantic and progress obligations. Programmatic strategy
synthesis~\cite{DBLP:journals/pacmpl/BatzBKW24} instead resolves
nondeterminism to obtain a selected program; it addresses a different
quantification over choices.
